\subsection{Additional Experimental Details}

\paragraph{Hyperparameters.} The narrowed benchmark runner records
the reader model, provider base URL, method name, temperature, top-$k$,
maximum materialised context length, sample mode, sample seed, and
sample size in every JSON report. Unless a live run overrides them,
the default values are temperature $0$, top-$k=6$, and
$120{,}000$ maximum context characters.

\paragraph{Infrastructure.} Offline harness validation was run locally
through \texttt{uv} and does not require PostgreSQL, Redis, GPU access,
or LLM credentials. Live Structure-service runs should additionally
record PostgreSQL, Redis, worker count, hardware, and provider network
conditions in the JSON artifact directory.

\paragraph{Reproducibility.} We release configuration files, seeds,
and the event logs produced by each run. Because the audit log is
deterministic, any reviewer can re-execute a trajectory against a
snapshot of the context store and obtain the same prompts.

\subsection{Context Tool Reference}

\begin{table}[h]
    \centering
    \small
    \begin{tabular}{ll}
        \toprule
        Tool & Purpose \\
        \midrule
        \texttt{glance}  & Scan one-line summaries under a prefix \\
        \texttt{list}    & Enumerate entries under a path \\
        \texttt{tree}    & Hierarchical view of a subtree \\
        \texttt{glob}    & Pattern match with \texttt{*}/\texttt{**} \\
        \texttt{read}    & Full \textsc{detail} for a single path \\
        \texttt{search}  & Semantic top-$k$ over embeddings \\
        \texttt{create}  & Insert a new entry \\
        \texttt{update}  & Modify content or metadata \\
        \texttt{delete}  & Remove an entry \\
        \texttt{rate}    & Rate a consulted source (\S\ref{sec:component_rating}) \\
        \bottomrule
    \end{tabular}
    \caption{The ten first-class context tools exposed to the agent.}
    \label{tab:tools}
\end{table}

\subsection{Default Event-GC Policy}
\label{sec:appendix_gc_policy}

Table~\ref{tab:ttl} lists the default TTL policy $\Pi$ used by the
event-level garbage collector (\S\ref{sec:component_gc}). The
\emph{pin} flag forces an event to remain visible for the life of the
run; otherwise the event expires once the decay accumulated from newer
events exceeds its TTL. Under the default decay rules, every newer
event adds $1$; any event ages \texttt{agent.token} and
\texttt{agent.heartbeat} by $4$ and \texttt{context.using} by $2$; an
\texttt{agent.message} ages \texttt{agent.thinking} by $3$ and a
\texttt{tool.call} ages it by $2$; a \texttt{tool.result} or
\texttt{tool.error} ages the matching \texttt{tool.call} and
\texttt{tool.pending} by $3$. A recency floor keeps the most recent
$100$ events per run ($500$ per workspace) regardless of type, and
event types not listed default to a TTL of $6$. Executors may override
any row per run.

\begin{table}[h]
    \centering
    \small
    \begin{tabular}{lcl}
        \toprule
        Event type & Steps & Notes \\
        \midrule
        \texttt{user.message}      & \textsc{pin} & Conversation turns \\
        \texttt{agent.message}     & \textsc{pin} & Agent replies \\
        \texttt{artifact.create/update/version} & \textsc{pin} & Durable outputs \\
        \texttt{task.complete}     & \textsc{pin} & Milestones \\
        \texttt{run.created}       & \textsc{pin} & Lifecycle anchor \\
        \texttt{run.completed/failed/cancelled} & \textsc{pin} & Terminal states \\
        \texttt{system.error}      & \textsc{pin} & Failure evidence \\
        \midrule
        \texttt{tool.call}         & 3 & Short plan--act window \\
        \texttt{tool.result}       & 3 & Matches call \\
        \texttt{tool.error}        & 6 & Keep for recovery \\
        \texttt{tool.pending}      & 1 & Superseded by result \\
        \texttt{tool.client.request} & 2 & Client round-trip \\
        \texttt{agent.thinking}    & 2 & Transient reasoning \\
        \texttt{agent.plan.step}   & 4 & Mid-run plans \\
        \texttt{agent.query}       & 2 & Retrieval intents \\
        \texttt{context.using}     & 1 & Visit marker \\
        \texttt{context.rated}     & 4 & Rating provenance \\
        \texttt{task.create}       & 10 & Workspace churn \\
        \texttt{task.update/delete/assign} & 6 & Workspace churn \\
        \texttt{run.state.change}  & 8 & Transition trail \\
        \texttt{system.notification} & 4 & Ambient signal \\
        \midrule
        \texttt{agent.token}       & 0 & Dropped from prompt \\
        \texttt{agent.heartbeat}   & 0 & Dropped from prompt \\
        \bottomrule
    \end{tabular}
    \caption{Default time-to-live policy for events in the agent's
    working context. Events dropped from the prompt remain in the
    persistent audit log.}
    \label{tab:ttl}
\end{table}

\subsection{Additional Qualitative Examples}

Annotated trajectories are included only when their live per-case JSONL
artifacts are present. Each example should link the final answer to
the selected chunks, evidence status, and token ledger used by the
runner.

\subsection{Prompts and Templates}

The live runner serialises the prompt inputs through method-specific
case JSONL files under the benchmark output directory. A submission
snapshot should archive those files together with the JSON summary so
the full prompting surface is reproducible.
